\section{Introducere}
\label{sec:introduction}

Jocurile de curse auto reprezintă unul dintre cele mai populare genuri din industria jocurilor video, combinând elemente de fizică, grafică 3D și inteligență artificială într-o experiență interactivă complexă \cite{rabin2009introduction}. Simularea realistă a dinamicii vehiculelor constituie o provocare inginerească semnificativă, deoarece necesită modelarea interacțiunii dintre roți, suspensie, motor și suprafața drumului \cite{pacejka2012tire}.

Proiectul prezentat în această lucrare constă în dezvoltarea unui joc de curse auto 3D utilizând motorul Unity (versiunea 6000.4.2f1), cu scopul de a demonstra aplicarea practică a conceptelor din domeniul Proiectării Avansate a Sistemelor (PAS). Obiectivele principale ale proiectului sunt:

\begin{itemize}
    \item \textbf{Fizica vehiculului jucătorului:} Implementarea unui model de conducere realist bazat pe componenta \texttt{WheelCollider} din Unity, cu suport pentru tracțiune integrală, frânare, coasting (rulare liberă), control al tracțiunii și asistență la aderență laterală.
    
    \item \textbf{Inteligența artificială a oponenților:} Proiectarea unui sistem de navigare bazat pe waypoint-uri (puncte de referință), care permite vehiculelor controlate de calculator să parcurgă autonom traseul de curse, cu adaptare dinamică a vitezei.
    
    \item \textbf{Definirea traseului prin waypoint-uri:} Construcția traseului de curse folosind un sistem de waypoint-uri configurabile direct în editorul Unity, care servește atât pentru navigarea AI, cât și pentru delimitarea traiectoriei circuitului.

    \item \textbf{Sisteme de joc:} Implementarea sistemelor auxiliare precum numărătoarea inversă la start, contorizarea tururilor, meniul de pauză și ecranul de sfârșit al cursei.
\end{itemize}

Lucrarea este organizată după cum urmează: Secțiunea~\ref{sec:related_work} prezintă lucrări conexe, Secțiunea~\ref{sec:methodology} detaliază arhitectura și implementarea, Secțiunea~\ref{sec:results} prezintă rezultatele obținute, iar Secțiunile~\ref{sec:future_work} și~\ref{sec:conclusion} discută direcții viitoare și concluzii.
